% =====================================================================================
% Topic 40 -- Future Generation Computer Systems (Elsevier), Full Length Article.
% elsarticle, two column, single spaced (FGCS Guide for Authors; see P7/FORMAT_FGCS_40.md).
% Build: bash code/paper/build_fgcs.sh [--final]
% Every number comes from numbers.tex / prereg_table.tex / sizing_table.tex, generated by
% code/paper/gen_numbers.py from results/*.json (shared with the IoT-J build). refs.bib is generated by
% code/paper/make_bib.py --style fgcs from live Crossref/DataCite metadata.
% =====================================================================================
\documentclass[final,5p,times,twocolumn]{elsarticle}
\usepackage[T1]{fontenc}
\usepackage{lmodern}
\usepackage{amsmath,amssymb,amsfonts,amsthm}
\usepackage{graphicx}
\usepackage{booktabs}
\usepackage{array}
\usepackage{url}
\usepackage[hidelinks]{hyperref}
\journal{Future Generation Computer Systems}
\input{numbers.tex}

\newtheorem{theorem}{Theorem}
\newtheorem{lemma}{Lemma}
\newtheorem{proposition}{Proposition}
\newtheorem{corollary}{Corollary}
\theoremstyle{definition}
\newtheorem{definition}{Definition}
\theoremstyle{remark}
\newtheorem{remark}{Remark}
\newcommand{\eps}{\varepsilon}
\newcommand{\Psif}{\Psi}
\newcommand{\lt}{\tilde{\ell}}
\newcommand{\E}{\mathbb{E}}

\begin{document}
\begin{frontmatter}
\title{Miss Floors and Storage Sizing for Energy-Harvesting Sensing Nodes Under Heavy-Tailed Events}
\author[fsb]{Dat Lam Quoc\corref{cor1}}
\ead{lamquocdat@gmail.com}
\cortext[cor1]{Corresponding author. ORCID: 0009-0004-5432-9343.}
\affiliation[fsb]{organization={FPT School of Business and Technology, FPT University},
  city={Ho Chi Minh City}, country={Viet Nam}}

\begin{abstract}
Energy-harvesting sensing nodes at the edge of IoT infrastructure run a detector on a duty cycle and miss events between activations and during outages when storage runs empty. We turn both into design-time certificates. For any event-independent schedule, the probability of missing an event is at least a closed-form floor that depends only on the mean activation rate and on the distribution of event durations; the floor is attained exactly when consecutive activations are spaced as evenly as whole slots allow. Coupling this floor with the storage dynamics gives an energy floor for a given storage capacity and a sandwich for the smallest capacity that meets a miss target. Together they form a sizing procedure that needs only a harvest trace, the node's power figures and event durations. We test the certificates on \nTraceCells{} harvest cells built from satellite irradiance and indoor light traces, with \nEv{} events from \nVid{} public videos, under a pre-registered train/test protocol. None of the \nFloorChecks{} floor checks is violated. Periodic operation needs \tsInvLo--\tsInvHi{} times more storage than the miss-optimal schedule at a given miss target, and average-day sizing underestimates the needed storage by a median factor of \avgHCMC{} (Ho Chi Minh City) to \avgIn{} (indoor). Lowering the frame rate when storage runs low does not reduce misses relative to plain blackout (median difference \thrMedSigned{} over \nThr{} cases); \EabstractSentence{}
\end{abstract}

\begin{keyword}
Energy harvesting \sep IoT edge nodes \sep Duty cycling \sep Storage sizing \sep Event detection \sep Energy-neutral operation \sep Design-time certificates
\end{keyword}
\end{frontmatter}

\section{Introduction}\label{sec:intro}
Self-sustaining IoT infrastructure increasingly pushes inference to the edge: a camera or sensor node that harvests solar or indoor light runs an object detector locally instead of streaming raw data. Such a node cannot afford to run the detector on every frame. It invokes the detector on a subset of slots, and it stops completely when its storage runs empty. Both decisions cost missed events, and both are fixed at design time by two provisioning choices: the duty cycle and the size of the storage element (battery or supercapacitor).

Design rules for these choices come from two separate traditions. Energy-neutral operation~\cite{kansal2007,buchli2014sensys} and real-time scheduling with energy curves~\cite{moser2007rts} answer \emph{when storage never empties}. Event capture with rechargeable sensors~\cite{kar2006ton,jaggi2009wn,ren2014tpds,yau2010tosn} answers \emph{how much a duty cycle misses} under stochastic recharge models. Model-based dimensioning frameworks for solar IoT nodes~\cite{kuaban2025iwcmc} size harvester and storage from an energy balance. What a system designer who provisions a fleet of such nodes needs is the link: for a measured harvest trace and a measured event-duration law, how small can the miss probability be for a given storage capacity, and how much capacity does a miss target require? The answer should be a certificate that can be computed before deployment and checked against data, not a simulation of one particular controller.

This paper gives that link in closed form and checks it on traces. The contributions are:
\begin{enumerate}
\item \emph{A schedule-free miss floor.} For every oblivious schedule with mean activation $a$, the miss probability is at least $\eps^*(a)$ (Theorem~\ref{thm:floor}). The floor is an elementary consequence of the Palm--Khinchin gap decomposition~\cite{daley2003pp1,lin2013opre,yau2010tosn}, and we state it as a certificate that is computed directly from an empirical duration law. The part that is new to our knowledge is the tightness condition: equality for every duration law holds if and only if all inter-activation gaps lie in $\{g,g+1\}$.
\item \emph{An energy floor and a sizing sandwich.} For a node with finite storage $C$, an energy-induced miss floor (Theorem~\ref{thm:energy}) and a sandwich for the smallest capacity $C_{\min}(\eps)$ that meets a miss target $\eps$ (Theorem~\ref{thm:sizing}), using the storage condition of~\cite{moser2007rts} and the offline-optimal energy-tunnel schedule of~\cite{tutuncuoglu2012twc}.
\item \emph{A sizing procedure for deployments} (Section~\ref{sec:procedure}) that needs only a harvest trace, the node's active/idle power ratio and leakage, and a sample of event durations, and that brackets what any event-independent controller can achieve.
\item \emph{A pre-registered trace-driven evaluation and an open artifact.} \nPass{} of \nPred{} frozen predictions pass on held-out data and two fail in magnitude but not in sign. The evaluation quantifies how far three common sizing rules are off, and tests spreading versus harvest-following schedules, cue-based gates, and throttling versus blackout. Code, pre-registration and derived data are released (see the Data availability statement).
\end{enumerate}
The node model is dimensionless (energies in units of the idle-slot energy, $r$ = active/idle ratio), so the results apply to any platform once $r$ and the leakage are known. Absolute numbers use one platform with published idle and active measurements, a microcontroller camera node~\cite{moosmann2025tinyissimo} ($r\approx14$); all other results sweep $r\in[1.1,300]$.

Section~\ref{sec:related} reviews related work. Section~\ref{sec:model} states the model, Sections~\ref{sec:floor}--\ref{sec:gate} give the certificates with proofs, Section~\ref{sec:eval} evaluates them on traces, Section~\ref{sec:procedure} turns them into a sizing procedure, and Section~\ref{sec:threats} discusses threats to validity.

\section{Related Work}\label{sec:related}
% related_work.tex (FGCS build): hand-written prose, no measured numbers. Every statement about a cited work is
% from its full text or, for works marked a in Table 1, from its abstract (P7/DOSSIER_FGCS_40.md, P6/KUABAN_40.md,
% P7/LIT_SWEEP_40.md, P1/DOSSIER_IOTJ_40.md).
Table~\ref{tab:related} positions the paper against the closest works along the dimensions that matter for provisioning a harvesting sensing node. We group the literature into three gaps.

\subsection{Sizing harvesting nodes targets continuity, not a miss rate}
The storage condition for a node that never empties its storage is classical: Moser et al.~\cite{moser2007rts} give it for energy curves, and Kansal et al.~\cite{kansal2007} give the energy-neutral sizing condition for a $(\rho,\sigma)$ harvest envelope. Prediction-based duty cycling~\cite{buchli2014sensys,renner2014tosn} and capacitor-centred designs~\cite{jackson2019ipsn} address operation rather than the event-miss trade-off. Recent design practice keeps the same objective. Kuaban~\cite{kuaban2025iwcmc} dimensions the harvester and storage of a solar-powered IoT node to maintain energy balance, using time-varying models and simulation; it does not consider event miss or event durations. A companion journal paper sizes the panel from peak sun hours with a tolerance factor and a design margin and sizes the battery from an autonomy target~\cite{kuaban2025sensors}, and a model of energy packets dimensions the harvester so that a node runs for several years without running out of energy~\cite{kuaban2024fenrg}. Field designs choose supercapacitors by days of autonomy~\cite{gurung2025ijdsn}, adapt the duty cycle to solar forecasts so that operation is continuous~\cite{gerndt2025sensors}, or size batteryless nodes by an energy-conservation inequality under constant illumination~\cite{landivar2025pimrc}. In this journal, Prauzek et al.~\cite{prauzek2025fgcs} tune fuzzy energy-management controllers for harvesting IoT nodes on several years of solar data, and Chen et al.~\cite{chen2025fgcs} allocate harvested energy optimally over a finite horizon in a linear sensor network, evaluated on a real renewable-energy data set. None of these sizes storage for a target miss probability of events; our Theorem~\ref{thm:sizing} generalises the zero-violation case to a miss target $\eps>0$ and Section~\ref{sec:eval} measures how far the continuity-oriented rules are from it.

\subsection{Event-capture analysis assumes stochastic recharge and point events}
Kar et al.~\cite{kar2006ton} bound network-wide capture under stochastic recharge with the concavity argument of Corollary~\ref{cor:jensen}. Jaggi et al.~\cite{jaggi2009wn,jaggi2008tosn} and Ren et al.~\cite{ren2014tpds,ren2015p2p} derive activation policies for correlated or renewal events; their capacity results are asymptotic or algorithmic. Quality of monitoring~\cite{yau2010tosn,he2012tpds} and patrolling~\cite{lin2013opre} supply the gap decomposition without energy. Quickest change detection with randomly available energy~\cite{naha2024tsp} and learned wake-up policies for event reporting~\cite{ruizguirola2026iotj} also model the harvest as a random process. Closed forms with finite storage exist for point arrivals under stochastic recharge: the energy-outage probability with Bernoulli energy arrivals~\cite{ballestrin2026cl}, blocking and depletion probabilities of a finite fluid queue~\cite{deepak2026mcap}, service-outage probabilities in a Poisson energy-packet model~\cite{kuaban2024fenrg}, and energy-depletion bounds under storage constraints~\cite{yan2024iotj}. We know of no closed-form finite-capacity floor for events with duration (Table~\ref{tab:related}; for the works marked $^{a}$ only the abstract was available to us).

\subsection{Offline-optimal harvesting schedules optimise throughput}
Offline-optimal transmission with a finite battery~\cite{tutuncuoglu2012twc,yang2012tcom,ozel2011jsac} and stochastic policies~\cite{sharma2010twc,srivastava2013ton} maximise throughput or utility, not the miss of events. We use the energy-tunnel construction of~\cite{tutuncuoglu2012twc} as the upper bound $C_{\mathrm{TS}}$ in Theorem~\ref{thm:sizing}, and recent IoT work bounds age of information~\cite{lei2025iotj,botirov2026iotj} under stochastic models.

\subsection{Intermittent systems and edge infrastructure are evaluated, not certified}
Batteryless cameras and intermittent platforms~\cite{desai2022camaroptera,hester2017mayfly,lee2026nowornever} show the hardware side. REACT~\cite{williams2024react} argues that a fixed energy buffer chosen at design time is a poor fit and that buffer dynamics are hard to predict at design time; it is evaluated on the same EnHANTs light traces we use. Flexible storage~\cite{alsubhi2024stash} and emulation platforms~\cite{liedtke2026estacker} evaluate coverage and performance over recorded solar and event traces, and self-sustainability studies use measured indoor-light traces~\cite{mayer2024iotj,kim2024iotj}. At infrastructure scale, Landr\'e et al.~\cite{landre2025fgcs} compute the minimum power demand for a renewable-powered datacenter under a quality-of-service constraint, and Kazem et al.~\cite{kazem2026fgcs} consolidate solar- and battery-powered edge clusters. These works evaluate a design on traces; our certificates bound every event-independent design before it is built, and the traces are used to check the bounds.

\begin{table*}[!t]\centering
\caption{Positioning (\checkmark: yes; $\circ$: partly; --: no). Read from full text unless marked $^{a}$ (abstract only; for these rows -- means not stated in the abstract).}\label{tab:related}
\footnotesize\setlength{\tabcolsep}{3pt}
\begin{tabular}{@{}lccccccl@{}}\toprule
Work & Energy & Events with & Finite & Closed-form & Sizing for & Real & Objective\\
 & model & duration & $C$ & bound & a target & traces & \\\midrule
Moser et al.~\cite{moser2007rts} & det.\ curve & -- & \checkmark & \checkmark & $\eps=0$ only & \checkmark & deadlines\\
Kansal et al.~\cite{kansal2007} & $(\rho,\sigma)$ envelope & -- & \checkmark & \checkmark & $\eps=0$ only & \checkmark & neutrality\\
Buchli et al.~\cite{buchli2014sensys} & daily prediction & -- & \checkmark & -- & \checkmark & \checkmark & continuity\\
Kuaban~\cite{kuaban2025iwcmc}$^{a}$ & time-varying model & -- & \checkmark & -- & $\circ$ & -- & energy balance\\
Prauzek et al.~\cite{prauzek2025fgcs}$^{a}$ & measured & -- & -- & -- & -- & \checkmark & transmissions\\
Chen et al.~\cite{chen2025fgcs}$^{a}$ & measured & -- & -- & $\circ$ & -- & \checkmark & throughput\\
Kar et al.~\cite{kar2006ton} & stochastic & -- & $\circ$ & \checkmark & -- & -- & capture (network)\\
Jaggi et al.~\cite{jaggi2009wn}$^{a}$ & stochastic & $\circ$ & asymptotic & $\circ$ & -- & -- & event capture\\
Ren and Cheng~\cite{ren2015p2p}$^{a}$ & stochastic & \checkmark & \checkmark & -- & -- & -- & event capture\\
Yau et al.~\cite{yau2010tosn} & none & \checkmark & -- & \checkmark & -- & -- & monitoring\\
Ballestrin and M\"uller~\cite{ballestrin2026cl}$^{a}$ & stochastic & -- & \checkmark & \checkmark & -- & -- & energy outage\\
P and Deepak~\cite{deepak2026mcap}$^{a}$ & stochastic & -- & \checkmark & \checkmark & -- & -- & blocking, depletion\\
Yan et al.~\cite{yan2024iotj}$^{a}$ & stochastic & -- & \checkmark & \checkmark & $\circ$ & -- & throughput\\
Tutuncuoglu and Yener~\cite{tutuncuoglu2012twc} & deterministic & -- & \checkmark & offline optimum & -- & -- & throughput\\
REACT~\cite{williams2024react} & measured & -- & \checkmark & -- & -- & \checkmark & buffer adaptation\\
Mayer et al.~\cite{mayer2024iotj} & measured & $\circ$ & \checkmark & -- & -- & \checkmark & sustainability\\
This work & det.\ trace & \checkmark & \checkmark & \checkmark & \checkmark & \checkmark & event miss\\\bottomrule
\end{tabular}
\end{table*}


\section{System Model}\label{sec:model}
\emph{Slots and energy.} Time is divided into slots of one frame period (100\,ms at 10\,fps) on a cyclic horizon of $T$ slots. In slot $t$ the node draws $e_t=1+\lt+x_t(r-1)$, where $x_t\in\{0,1\}$ indicates a detector activation and $\lt\ge0$ is a constant leakage. Energies are in units of the idle-slot energy $e_{\mathrm{idle}}$. The harvest is $h_t\ge0$, the capacity is $C$, and the storage evolves as $B_t=\min(C,B_{t-1}+h_t-e_t)$. If a slot cannot be paid, the node is in \emph{outage}: it skips activations and restarts when $B\ge0.05C$.

\emph{Policies.} A schedule $x\in\{0,1\}^T$ has mean activation $a(x)=T^{-1}\sum_t x_t$. It is \emph{oblivious} if it does not depend on the event process. It may depend on time, on the harvest trace (even offline) and on its own storage. A \emph{gate} pays a cue cost $c$ per slot and activates according to frame content.

\emph{Events.} An event has onset $s$ and duration $D\in\{1,2,\dots\}$ slots. It is caught if some activation falls in $\{s,\dots,s+D-1\}$; $\mathrm{miss}(x)$ is the probability that it is not.

\emph{Assumption A1.} Onsets are uniform on the horizon, $D\sim F_D$, and $(s,D,x)$ are independent.

\emph{Notation.} $\phi(g)=\E[\min(g,D)]$ and $\Psif(a)=\E[\min(1,aD)]$ are concave. $\eps^*(a)=1-\Psif(a)$ is convex and non-increasing.

\section{A Miss Floor for Oblivious Schedules}\label{sec:floor}
\begin{lemma}[Gap decomposition; restated]\label{lem:gap}
Let a deterministic schedule have activations at cyclic positions $\tau_1<\dots<\tau_n$ with gaps $g_i=\tau_{i+1}-\tau_i$. Under A1, $1-\mathrm{miss}(x)=T^{-1}\sum_{i=1}^n\phi(g_i)$.
\end{lemma}
\begin{proof}
Onsets $s\in(\tau_i,\tau_{i+1}]$ are $g_i$ positions. For each, the first activation at or after $s$ is $\tau_{i+1}$, at distance $k\in\{0,\dots,g_i-1\}$, and the event is caught iff $D>k$. Summing $P(D>k)$ over $k$ gives $\phi(g_i)$; averaging over the uniform onset gives the claim.
\end{proof}
\noindent This is the discrete Palm--Khinchin form~\cite{daley2003pp1}. It is the per-gap cost of patrolling with random attack times~\cite{lin2013opre}, and for one gap value it is the periodic quality-of-monitoring formula~\cite[Thm.~4.2]{yau2010tosn}.

\begin{theorem}[Miss floor]\label{thm:floor}
(i) Every oblivious schedule with $\E[a(x)]=a$ has $\mathrm{miss}\ge\eps^*(a)$.
(ii) For deterministic $x$, equality holds iff for every $d$ with $P(D=d)>0$ either all gaps are $\le d$ or all gaps are $\ge d$. Hence equality holds for every $F_D$ iff all gaps lie in $\{g,g+1\}$ for some integer $g$. Periodic and interleaved schedules with gaps $\lfloor1/a\rfloor$ and $\lceil1/a\rceil$ attain the floor.
\end{theorem}
\begin{proof}
(i) For every $d$, $\sum_i\min(g_i,d)\le\min(\sum_ig_i,nd)=\min(T,nd)$. Taking expectations over $D$ and using Lemma~\ref{lem:gap}, $1-\mathrm{miss}=T^{-1}\E_D\sum_i\min(g_i,D)\le\E_D[\min(1,(n/T)D)]=\Psif(n/T)$. For a random oblivious schedule, condition on $x$ and apply Jensen's inequality to the concave $\Psif$. The slack is exact: $\mathrm{miss}(x)-\eps^*(a)=T^{-1}\E_D[\min(T,nD)-\sum_i\min(g_i,D)]\ge0$.
(ii) $\sum_i\min(g_i,d)=\min(T,nd)$ iff all $g_i\le d$ (both sides equal $T$) or all $g_i\ge d$ (both sides equal $nd$); if $g_i<d<g_j$ for some $i,j$, the left side is below both. Gaps in $\{g,g+1\}$ leave no integer strictly between two gaps. Conversely, two gaps that differ by at least 2 admit an integer $d$ strictly between them, and a point mass at $d$ breaks equality. Interleaved (Bresenham) schedules with gaps $\lfloor1/a\rfloor$ and $\lceil1/a\rceil$ have gaps in such a pair.
\end{proof}
\noindent Part (i) is an elementary certificate; its value is that it is computed from an empirical $F_D$ without knowing the schedule. Part (ii) shows that the floor is exact, not only a bound, for the schedules a designer would use. On our corpus, $P(D=1)=\pDone$ and \pDltEight{} of events are shorter than 8 frames. As a result $\eps^*(1/8)=\epsKeight$, even though periodic $K=8$ catches every event with $D\ge8$.

\begin{corollary}[Spread beats burst]\label{cor:jensen}
If a schedule runs at rate $a_j$ on blocks of length $n_j\gg D_{\max}$ with $\sum_jn_ja_j=aT$, then $1-\mathrm{miss}\approx\sum_j(n_j/T)\Psif(a_j)\le\Psif(a)$.
\end{corollary}
\begin{proof}
Events inside block $j$ are caught with probability $\Psif(a_j)$ by Theorem~\ref{thm:floor}(ii) applied to the block; only events that straddle a block boundary are affected, $O(\E[D])$ per boundary. Concavity of $\Psif$ gives the inequality.
\end{proof}
\noindent The Jensen gap is the exact cost of a non-constant rate. It is the same principle as~\cite[Thm.~1]{kar2006ton}, applied to a harvest-driven rate profile and to events with duration. Harvest-following, day-only operation and outages are all special cases.

\section{Energy Floor and Storage Sizing}\label{sec:energy}
\begin{lemma}[Storage never empties; restated]\label{lem:moser}
Starting full, $B_t\ge0$ for all $t$ iff $\sum_{u=s+1}^t(e_u-h_u)\le C$ for all $s<t$~\cite[Thm.~4,5]{moser2007rts}; leakage enters as in~\cite[Eq.~8]{kansal2007}. The minimal capacity is the maximum drawdown of the cumulative net income.
\end{lemma}
\begin{proof}
Without the lower clip, the recursion with $B_0=C$ has the closed form $B_t=\min_{0\le s\le t}(C+H(s,t]-E(s,t])$, a reflected walk. Non-negativity of every term is the window condition. On a cyclic horizon in steady state the minimal $C$ is the maximum drawdown over two concatenated cycles, and it is infinite when the cycle's net income is negative.
\end{proof}

Let $\nu_t$ be the least storage at the end of slot $t$ from which idle-only operation survives to the horizon: $\nu_T=0$, $\nu_{t-1}=[\nu_t+1+\lt-h_t]^+$. If $\nu_t>C$ for some $t$, no schedule avoids an outage.

\begin{lemma}[Largest activation budget]\label{lem:greedy}
Starting full, the greedy schedule that spends $y_t=\min\{r-1,\,B_{t-1}+h_t-1-\lt-\nu_t\}^+$ on activation in every slot maximises $\sum_ty_t$ over all schedules that never empty storage.
\end{lemma}
\begin{proof}
Let $P$ be any feasible schedule. By induction, $B^G_t\le B^P_t$: either the greedy stopped at $\nu_t\le B^P_t$, or it spent the maximum while starting from no more storage. Hence the greedy never spills more than $P$. With $Z_t$ the cumulative spill plus $B_t$, $Z_t=Z_{t-1}+h_t-1-\lt-y_t$ gives $Z^G_T\le Z^P_T$, so $\sum_ty^G_t\ge\sum_ty^P_t$.
\end{proof}
\begin{remark}[Fluid relaxation]
Lemma~\ref{lem:greedy} lets the activation energy $y_t$ take any value in $[0,r-1]$, whereas a real node spends $x_t(r-1)$ with $x_t\in\{0,1\}$. Every binary schedule that never empties storage is a feasible point of this relaxation, so the relaxed maximum is an upper bound on the mean activation of every binary no-outage schedule. The upper bound is all that Theorem~\ref{thm:energy} uses.
\end{remark}
Let $A(C)$ be the mean activation of the greedy schedule. In cyclic steady state the energy balance adds $a\le a_{\mathrm{bal}}=(\bar h-1-\lt)/(r-1)$; let $A_{\mathrm{cyc}}(C)=\min(A(C),a_{\mathrm{bal}})$. The single-pass $A(C)$ can exceed $a_{\mathrm{bal}}$ when $C$ is large relative to the season's demand (it starts full), which is why the cyclic bound takes the minimum.

\begin{theorem}[Energy-induced miss floor]\label{thm:energy}
Every oblivious schedule that never empties storage on $(h,C)$ has $\mathrm{miss}\ge\eps^*(A_{\mathrm{cyc}}(C))$, even with offline knowledge of $h$.
\end{theorem}
\begin{proof}
Its mean activation is at most $A_{\mathrm{cyc}}(C)$ (Lemma~\ref{lem:greedy}, the remark, and the cyclic energy balance), $\eps^*$ is non-increasing, and Theorem~\ref{thm:floor} applies.
\end{proof}

\begin{proposition}[Blackout]\label{prop:blackout}
An oblivious schedule with outage fraction $\pi$, $N$ outage episodes and on-time activation $a_{\mathrm{on}}$ has $\mathrm{miss}\ge\pi+(1-\pi)\eps^*(a_{\mathrm{on}})-N\E[D]/T$.
\end{proposition}
\begin{proof}
Events whose span lies inside an outage are missed. Concatenating the on-slots into one cyclic horizon and applying Theorem~\ref{thm:floor} bounds the rest; only events that straddle an on/off boundary are affected by the concatenation, at most $\E[D]$ per boundary in expectation.
\end{proof}
\noindent At trace scale outages last hours and $\E[D]$ is a few seconds, so the correction is negligible. The term $\pi$ is the loss-of-power-supply probability of storage sizing~\cite{borowy1996lpsp}. Only the second term is new. By convexity, $\pi+(1-\pi)\eps^*(a_{\mathrm{on}})\ge\eps^*((1-\pi)a_{\mathrm{on}})$: an outage is the most bursty rate profile.

For a miss target $\eps$ let $a_\eps=\min\{a:\eps^*(a)\le\eps\}$ and define $C_{\min}(\eps)$ as the least $C$ at which some oblivious no-outage schedule reaches $\mathrm{miss}\le\eps$. Further, $C_{\mathrm{LB}}(\eps)=\inf\{C:A_{\mathrm{cyc}}(C)\ge a_\eps\}$, $C_{\mathrm{per}}(\eps)$ is the capacity needed by a periodic schedule at $a_\eps$ (Lemma~\ref{lem:moser}), and $C_{\mathrm{TS}}(\eps)$ is the least $C$ at which the taut-string schedule reaches $\eps$.

\begin{theorem}[Sizing sandwich]\label{thm:sizing}
$C_{\mathrm{LB}}(\eps)\le C_{\min}(\eps)\le\min\{C_{\mathrm{TS}}(\eps),C_{\mathrm{per}}(\eps)\}$. If some miss-optimal no-outage cyclic schedule is full at the taut-string anchor and stays below rate 1, then $C_{\min}(\eps)=C_{\mathrm{TS}}(\eps)$.
\end{theorem}
\begin{proof}
(a) A schedule with miss $\le\eps$ has $a\ge a_\eps$ by Theorem~\ref{thm:floor} and $a\le A_{\mathrm{cyc}}(C)$ by Theorem~\ref{thm:energy}; $A_{\mathrm{cyc}}$ is non-decreasing in $C$, which gives the left inequality.
(b) The taut string (by construction) and periodic operation at $a_\eps$ (Lemma~\ref{lem:moser}) are both feasible no-outage schedules, which gives the right inequality.
(c) Without the rate cap, a spilling schedule is dominated by a non-spilling one that consumes the would-be spill: the storage path is the same while storage is not full, the rates are pointwise higher, and $\Psif$ is increasing. A non-spilling cyclic schedule that is full at the anchor $t^*$ is a path of the cumulative activation energy $Y$ in the tube $S_j\le Y_j\le C+S_j$, with $S_j$ the cumulative idle-only surplus. Among such paths the taut (shortest) string maximises $\sum_j\Psif(\Delta Y_j)$. For strictly concave rate functions this is~\cite[Thm.~1]{tutuncuoglu2012twc}, whose energy-tunnel Algorithm~A1 uses only the arrivals and the storage bound, not the function (see also~\cite{yang2012tcom}). For the piecewise-linear $\Psif$, apply it to $\Psif+\delta\sqrt{\cdot}$ and let $\delta\to0$: the optimal value is continuous in $\delta$ and the path does not depend on it.
\end{proof}
\noindent Periodic operation is miss-optimal at a given rate (Theorem~\ref{thm:floor}), but not storage-optimal at a given miss. The taut string keeps the rate as flat as the energy tube allows. It runs faster while the sun shines and slower before a deficit. In the evaluation $C_{\mathrm{TS}}$ is found by bisection with relative tolerance 0.01.

\section{Gates with a Cue Cost}\label{sec:gate}
A gate is not oblivious, so Theorems~\ref{thm:floor}--\ref{thm:sizing} do not bound its miss. Its energy, however, is bounded below.
\begin{proposition}[Gate boundary]\label{prop:gate}
Suppose a gate activates at long-run rate $a_g=\rho\,a_{\mathrm{per}}$, spread over time, and pays $c$ per slot. At the capacity $C^*_{\mathrm{per}}$ that exactly sustains periodic $a_{\mathrm{per}}$, the gate empties storage iff $c>c^*(\rho)=(1-\rho)a_{\mathrm{per}}(r-1)$.
\end{proposition}
\begin{proof}
Relative to periodic operation, the gate changes the demand per slot by $c+(a_g-a_{\mathrm{per}})(r-1)$. The window condition of Lemma~\ref{lem:moser} with this extra constant demand fails at $C=C^*_{\mathrm{per}}$ iff the change is positive, since $C^*$ is strictly increasing in a constant demand. If the change exceeds the periodic slack per slot $\bar h-(1+\lt+a_{\mathrm{per}}(r-1))$, the gate empties storage for every finite $C$.
\end{proof}
\noindent The boundary does not depend on the site; the site enters only through the extra storage and outage. An ideal presence gate on our corpus would be active on \inEventFrac{} of frames, i.e.\ $\rho\approx\rhoPresence$ against $K=8$, so it is infeasible at $C^*_{\mathrm{per}}$ for every $c\ge0$. The simulated map agrees with the analytic boundary in \nGateAgree{} of \nGate{} cells (Fig.~\ref{fig:gate}).

\begin{figure}[!t]\centering
\includegraphics[width=\columnwidth]{figures/F4_gate_map_test.pdf}
\caption{Median outage of a gate run at $C^*_{\mathrm{per}}(8)$ against the rate ratio $\rho$ and the cue cost $c$ ($r=15$, $\lt=0.04$, $m=0.25$; TEST). The staircase is the analytic boundary $c^*(\rho)$ of Proposition~\ref{prop:gate} on cell edges; the dashed line marks an ideal presence gate.}\label{fig:gate}
\end{figure}

\section{Trace Evaluation}\label{sec:eval}
\subsection{Setup and pre-registration}\label{sec:setup}
\emph{Harvest.} Satellite global horizontal irradiance for 2025 at native resolution (10--30\,min)~\cite{zippenfeld2023openmeteo}: Ho Chi Minh City (Himawari-9), Madrid and Uppsala (MSG and SARAH-3~\cite{pfeifroth2023sarah3}), 4 seasons each. Indoor irradiance from EnHANTS setups B--F~\cite{enhants2011dataset,gorlatova2013tmc}. This gives \nTraceCells{} cells. Traces are held on 1-min sub-steps (zero-order hold; linear interpolation changes $C^*$ by at most \interpMax). The panel is sized so that the mean harvest is $(1+m)$ times the demand of periodic $K=8$, $m\in\{0.25,1\}$.

\emph{Events.} The duration law $F_D$ is rebuilt from detector outputs~\cite{ultralytics_yolo} on \nVid{} public videos (CDnet~2014~\cite{wang2014cdnet}, LASIESTA~\cite{cuevas2016lasiesta}, BMC~\cite{vacavant2013bmc}). An event is a maximal run of frames with at least one detection; there are \nEv{} events with $\E[D]=\ED$ frames. Only $F_D$ enters the results.

\emph{Grid and protocol.} The grid crosses $r\in\{1.1,\dots,300\}$ (8 values), $\lt\in\{0,0.01,0.04,0.1\}$, $m$, and 20 capacities per combination, with eleven policies (six periodic with $K\in\{2,\dots,64\}$, throttling, harvest-following, a Kansal-style EWMA duty cycle, taut string, greedy). Odd months and half of the videos (seeded) form TRAIN; even months and the other videos form TEST. Predictions with pass criteria were frozen (SHA-256 and a git tag) after TRAIN and before TEST. Every number below is from TEST. Misses are measured by replaying $10^5$ real events per run; bootstrap intervals resample videos ($B=1000$) or 7-day blocks of the trace. Sign tests use Holm's correction at family-wise $\alpha=0.05$. Table~\ref{tab:prereg} lists every prediction, its criterion and its outcome. The two failures are reported, not corrected.

\newcommand{\preregCritW}{0.49\textwidth}\newcommand{\preregObsW}{0.32\textwidth}
\input{prereg_table.tex}

\subsection{The floors hold and are tight}
Across \nTone{} policy--capacity runs, no replayed miss lies below the floor of Theorem~\ref{thm:floor} by more than a Bonferroni margin ($z=\zBonf$). The same holds for the \nPtwo{} blackout runs (Proposition~\ref{prop:blackout}) and the \nTtwo{} no-outage runs (Theorem~\ref{thm:energy}); the latter also have no analytic violation above $10^{-6}$. Fig.~\ref{fig:floor} shows that the periodic and taut-string runs lie on the floor. Every point above it is a policy with a variable rate or with outages.

\begin{figure}[!t]\centering
\includegraphics[width=\columnwidth]{figures/F2_miss_vs_floor_test.pdf}
\caption{Replayed miss against the applicable floor (Theorem~\ref{thm:floor}, \ref{thm:energy} or Proposition~\ref{prop:blackout}) for a random sample of TEST runs. \Ftwocaption}\label{fig:floor}
\end{figure}

\subsection{Spreading versus following the harvest}
At equal mean rate $a=1/8$, harvest-following needs \dCMinPct--\dCMaxPct{} less storage than periodic operation but misses \dmissMin--\dmissMax{} more events (significant in all \nBootCells{} bootstrap cells after Holm over \nHolm{} tests, largest $p=\pMaxBoot$). At equal miss target the comparison reverses: the taut string needs only \tsMedLo--\tsMedHi{} of the periodic capacity (medians per site group, $\eps\in\{0.5,0.6\}$, \nEps{} cells, ratio below 1 in all of them). Fig.~\ref{fig:frontier} shows both comparisons. The lower bound $C_{\mathrm{LB}}$ is within \lbtsGapPct{} of $C_{\mathrm{TS}}$ for $\eps\ge0.6$ (medians $\ge\lbtsHi$), and loose at $\eps=0.45$ (\lbtsLoMin--\lbtsLoMax).

\begin{figure}[!t]\centering
\includegraphics[width=\columnwidth]{figures/F3_frontier_spread_vs_hf_test.pdf}
\caption{Left: harvest-following against periodic at $a=1/8$ ($r=15$, $\lt=0.04$, $m=0.25$; one point per cell). Right: taut-string capacity relative to periodic at equal miss target (median per site).}\label{fig:frontier}
\end{figure}

\subsection{Sizing}
Table~\ref{tab:sizing} gives the storage needed by a class-A node ($r=14.3$, $e_{\mathrm{idle}}=0.248$\,mJ, $\lt=0.036$, $K=8$) over the full year. The baselines are average-day sizing, Buchli's daily-resolution deficit~\cite{buchli2014sensys}, Kansal's condition with the mean harvest rate~\cite{kansal2007}, and a 3-day autonomy rule. Across the full grid, average-day sizing underestimates $C^*$ by median factors of \avgHCMC{} (Ho Chi Minh City), \avgMadrid{} (Madrid), \avgUpp{} (Uppsala) and \avgIn{} (indoor), with maxima of \avgMaxOut{} outdoors and \avgMaxIn{} indoors. Daily resolution hides the night in the tropics (Buchli median \buHCMC{} in Ho Chi Minh City, \buLo--\buHi{} elsewhere). The 3-day rule always oversizes in Ho Chi Minh City (largest ratio \threeHCMCmax) but undersizes Uppsala by up to \threeUppMax{} times. Kansal's tight condition reproduces $C^*$ exactly, as it must; with the mean harvest rate it oversizes by \kmeanInvLo--\kmeanInvHi{} times. The sizing table is descriptive and was not pre-registered. The pre-registered TEST version (one month per season) failed its magnitude prediction (H7) because $C^*$ of a single month depends on that month's worst spell (supplementary material).

\input{sizing_table.tex}

The period itself is the cheapest lever in winter. Fig.~\ref{fig:ck} shows the storage need of periodic operation against $K$: doubling $K$ from 8 to 16 cuts the winter (DJF) need by a median factor of \ckWinter{} at the European sites, and by \ckHCMC{} in Ho Chi Minh City. The price is a miss of $\eps^*(1/16)=\epsKsixteen$ instead of \epsKeight.

\begin{figure}[!t]\centering
\includegraphics[width=\columnwidth]{figures/F1_frontier_K_Cstar_test.pdf}
\caption{Storage need of periodic operation against the period $K$ in days of demand ($r=15$, $\lt=0.04$, $m=0.25$; TEST months), winter (DJF) and summer (JJA).}\label{fig:ck}
\end{figure}

\subsection{Throttling, blackout and spending the spill}
Throttling lowers the activation rate when storage is low instead of letting the node black out. In the \nThr{} TEST cases where periodic $K=8$ has outages, even the best of five throttling variants per case (chosen with hindsight) stays within $\pm0.03$ of blackout in all cases (median \thrMedSigned). The common rule ``drop from $K=8$ to $K=64$ below 20\,\% storage'' misses more than blackout in \naiveSharePct{} of cases (median \naiveMedSigned). Outages last hours, so a throttled node still pays the idle energy and enters the same deficit.

\Esection\label{sec:explore}

\section{A Sizing Procedure for Deployments}\label{sec:procedure}
The results combine into a provisioning procedure that a fleet planner can run per site before deployment. It needs only a harvest trace, the node's $r$ and $\lt$, and a sample of event durations from the target scene.
\begin{enumerate}
\item Estimate $F_D$ and tabulate $\eps^*(a)$; the miss target fixes the activation rate $a_\eps$.
\item The energy balance bounds the reachable floor by $\eps^*(a_{\mathrm{bal}})$; a target below it needs a larger panel, not more storage.
\item Size storage with $C_{\mathrm{per}}(\eps)$ (Lemma~\ref{lem:moser}) for a constant-rate node, or with $C_{\mathrm{TS}}(\eps)$ for a node that knows its harvest; Theorem~\ref{thm:sizing} brackets what any event-independent controller can do, so a vendor controller that claims to beat $C_{\mathrm{LB}}(\eps)$ violates the model.
\item Size on the worst season and on more than one month of data, and report the interval, because $C^*$ is set by the single worst spell.
\item Do not rely on threshold throttling to soften outages; spend the spill with a predictive duty cycle, or add storage or panel (Section~\ref{sec:explore}).
\item Before adding a content-based gate, check its cue cost against $c^*(\rho)$ (Proposition~\ref{prop:gate}); a gate that is active on most event frames cannot save energy at the periodic capacity.
\end{enumerate}
\emph{Who uses it.} The procedure is meant for the planning stage of self-sustaining IoT deployments: choosing a supercapacitor or battery per site class, choosing the frame period of a detector firmware, and checking a vendor's claimed miss rate against a certificate. Every step is a closed-form computation or a maximum drawdown on a trace, so it runs in seconds per site and can be embedded in a deployment-planning tool; the released artifact (Data availability statement) contains a reference implementation.

\emph{Relation to companion work.} The event corpus and the affine energy model follow the author's earlier journal paper on frame skipping~\cite{deep3_jiot2026}. Here the corpus is rebuilt from public data by this paper's own pipeline and checked to be identical. No theorem of that paper is used.

\section{Threats to Validity}\label{sec:threats}
(i) Onsets are assumed uniform (A1). Real event rates are diurnal, which weights $\Psif$ by time of day and breaks the site independence of $\eps^*$; the corpus has no wall-clock timestamps, so we could not test this. (ii) The corpus is surveillance video; only $F_D$ is used, and the floor is computable for any other duration law. (iii) Satellite products disagree at high latitude: in Uppsala, $C^*$ from MSG and SARAH-3 differs by up to \uppProdDiffMax{} (median \uppProdDiffMed), so we report both. (iv) TEST has one month per season, which made the absolute sizing prediction fail. (v) $C^*$ is a maximum drawdown; block bootstrap intervals for such statistics can exclude the point estimate, so we report \blockNote. (vi) The taut string is an offline bound, not a deployable policy. (vii) Results are design-time certificates from traces and published power figures; no hardware was measured.

\section{Conclusion}\label{sec:conclusion}
Two closed-form quantities, the miss floor $\eps^*(a)$ and the energy-limited activation $A_{\mathrm{cyc}}(C)$, determine how often an energy-harvesting sensing node can catch events and how much storage it needs. On held-out traces the floors are never violated, periodic schedules attain them, and common sizing rules are off by large factors: average-day sizing by a median of up to \avgIn{} times indoors, and Kansal's mean-rate condition by \kmeanInvLo--\kmeanInvHi{} times in the other direction. For the designers of self-sustaining IoT infrastructure the practical message is a short procedure: size from the duration law and the worst spell of a measured trace, spread activations, and spend surplus energy rather than throttle.

\section*{CRediT authorship contribution statement}
\textbf{Dat Lam Quoc:} Conceptualization, Methodology, Formal analysis, Software, Validation, Investigation, Data curation, Visualization, Writing -- original draft, Writing -- review \& editing.

\section*{Declaration of competing interest}
The author declares that there are no known competing financial interests or personal relationships that could have appeared to influence the work reported in this paper.

\section*{Funding}
This research did not receive any specific grant from funding agencies in the public, commercial, or not-for-profit sectors.

\section*{Declaration of generative AI and AI-assisted technologies in the manuscript preparation process}
During the preparation of this work, the author used Claude (Anthropic) in order to develop code, search the literature and edit the language. After using this tool, the author reviewed and edited the content as needed and takes full responsibility for the content of the published article.

\section*{Data availability}
Code (MIT licence), the pre-registration and the derived data (CC BY 4.0) that generate every number, table and figure of this paper are available at \repoURL{} (tag \texttt{v1.1.0-fgcs}). The released event corpus contains onsets and durations only. The input data are third-party public data and are not redistributed: satellite irradiance from the Open-Meteo satellite API (CC BY 4.0)~\cite{zippenfeld2023openmeteo,pfeifroth2023sarah3}, indoor irradiance from the EnHANTS data set (CC BY 4.0)~\cite{enhants2011dataset}, and the CDnet~2014, LASIESTA and BMC videos~\cite{wang2014cdnet,cuevas2016lasiesta,vacavant2013bmc}.

\bibliographystyle{elsarticle-num}
\bibliography{refs}

\section*{Vitae}
\noindent
\begin{minipage}[t]{19mm}
\vspace{0pt}\includegraphics[width=19mm]{author_photo.jpg}
\end{minipage}\hfill
\begin{minipage}[t]{\dimexpr\linewidth-23mm\relax}
\vspace{0pt}\textbf{Dat Lam Quoc} is a PhD researcher at the FPT School of Business and Technology, FPT University, Ho Chi Minh City, Viet Nam. The author's research is on certified resource management for edge perception: duty cycling, energy harvesting and worst-case guarantees for event detection on constrained IoT nodes, checked against public traces with pre-registered protocols.
\end{minipage}
\end{document}
